%% kapa: exam=1 round=23 year=2012 period=1 form=A subject=civil_law subject_ko=민법 seq=28 number=28
%% source: Q-Net 감정평가사 기출문제 (cst00309) artlSeq=5052578 file=raw/1008/5052578/2012년 제23회 감정평가사 제1차시험 기출문제(1교시 A형).pdf
%% conversion: kordoc 4.24.51 -> tools/gen1.py (auto); review=proofread-sonnet; flags=none
\begin{question}{28}
등기부상 甲 소유의 X 토지를 乙이 소유의 의사로 평온ㆍ공연하게 1989. 2. 1.부터 2012. 7. 1.까지 점유하고 있다. 다음 설명 중 옳지 않은 것은? (다툼이 있으면 판례에 의함)
\begin{choices}
  \item 2009. 3. 2. 甲이 丙에게 소유권이전등기 후 다시 20년이 경과하더라도 乙은 등기명의자인 丙에게 취득시효를 주장할 수 없다.
  \item 2012. 7. 1. 甲이 X 토지의 반환을 청구하는 경우, 乙은 등기 없이도 甲의 반환청구에 대항할 수 있다.
  \item 甲은 乙에게 1989. 2. 1.부터 취득시효 완성시까지의 임료 상당액을 반환청구할 수 없다.
  \item 2010. 3. 5. 甲이 丁에게 X 토지에 대한 저당권설정등기를 해준 경우, 특별한 사정이 없는 한 저당권설정등기는 유효하다.
  \item 1995. 5. 15. 甲이 戊에게 소유권이전등기를 해준 경우, 2012. 7. 1. 乙은 戊에 대하여 소유권이전등기를 청구할 수 있다.
\end{choices}
\end{question}
\answer{1}
